





%% =====================================================================
\section{Complexity Analysis}
\label{sec:analysis}

In this section, we analyze the time and space complexity of the projection-based filter, characterize its interaction with the NN-Descent algorithm, and derive the conditions under which filtering yields a net computational benefit.

\subsection{Preprocessing Cost}
\label{sec:preprocessing}

The one-time preprocessing phase consists of two steps:

\paragraph{Projection vector generation.}
Generating $m$ random Gaussian vectors of dimension $d$ costs $O(m \cdot d)$.
This is negligible relative to the projection step.

\paragraph{Data projection.}
Computing the projections $P(x_i) \in \mathbb{R}^m$ for all $n$ data points requires $n \times m$ inner products, each of cost $O(d)$, yielding a total preprocessing cost of:
\begin{equation}
    T_{\text{preprocess}} = O(n \cdot m \cdot d)
    \label{eq:preprocess}
\end{equation}
For our experimental setting ($n = 10^6$, $m = 32$, $d = 960$), this amounts to approximately $3 \times 10^{10}$ floating-point operations.
In practice, this step takes a few seconds and is amortized over all subsequent iterations.

\subsection{Per-Iteration Cost With Filtering}
\label{sec:per_iteration}

In standard NN-Descent, each iteration generates $C$ candidate pairs and computes an exact distance for each, at a cost of $O(C \cdot d)$ per iteration.
With projection-based filtering, the per-iteration cost becomes:
\begin{equation}
    T_{\text{iter}} = \underbrace{C \cdot O(m)}_{\text{filter check (all pairs)}} + \underbrace{(1 - f) \cdot C \cdot O(d)}_{\text{distance computation (surviving pairs)}}
    \label{eq:iter_cost}
\end{equation}
where $f \in [0, 1]$ is the \emph{filter rate}---the fraction of candidate pairs that are filtered (skipped) per iteration.
Rearranging, the cost ratio relative to unfiltered NN-Descent is:
\begin{equation}
    \frac{T_{\text{filtered}}}{T_{\text{unfiltered}}} = \frac{m}{d} + (1 - f) = 1 - f + \frac{m}{d}
    \label{eq:cost_ratio}
\end{equation}

\subsection{Break-Even Condition}
\label{sec:breakeven}

Filtering provides a net speedup when $T_{\text{filtered}} < T_{\text{unfiltered}}$, i.e., when:
\begin{equation}
    1 - f + \frac{m}{d} < 1 \quad \Longrightarrow \quad f > \frac{m}{d}
    \label{eq:breakeven}
\end{equation}
This gives a clean criterion: \textbf{the filter rate must exceed the ratio of projection dimensionality to data dimensionality}.

\begin{itemize}
    \item \textbf{GIST ($d = 960$, $m = 32$):} Break-even at $f > 3.3\%$. Our observed filter rates range from 15--25\% (Table~\ref{tab:gist1m_results}), well above the threshold.
    \item \textbf{SIFT ($d = 128$, $m = 32$):} Break-even at $f > 25\%$. The observed filter rates are comparable but the margin is thin, explaining the marginal speedup on this dataset.
\end{itemize}

This analysis also explains \emph{why} the filter is most beneficial for high-dimensional data: as $d$ grows, the break-even threshold $m/d$ decreases, making even modest filter rates profitable.

\subsection{Memory Overhead}
\label{sec:memory}

The filter requires storing the precomputed projections for all $n$ points:
\begin{equation}
    M_{\text{proj}} = n \cdot m \cdot \text{sizeof(float)} = n \cdot m \cdot 4 \;\text{bytes}
    \label{eq:memory}
\end{equation}
For $n = 10^6$ and $m = 32$, this amounts to 128~MB.
For comparison, the original dataset occupies $n \cdot d \cdot 4 = 3{,}840$~MB for GIST~1M.
The projection storage is therefore approximately $m / d = 3.3\%$ of the dataset size---a modest overhead.

Additionally, the filter requires $O(m \cdot d)$ storage for the $m$ projection vectors themselves, which is negligible.

\subsection{Adaptive Filter Behavior}
\label{sec:adaptive}

An important property of the projection-based filter is that it is \emph{self-adaptive} across iterations without any explicit parameter adjustment.
The filter condition (Equation~\ref{eq:proj_filter}) depends on $d_k(u)$---the distance to each point's current farthest neighbor---which \emph{decreases} as the graph improves.
As $d_k(u)$ shrinks, the threshold $t^2 \cdot d_k(u)^2$ tightens, making the filter condition easier to satisfy and thus filtering more aggressively.

This behavior is desirable: in early iterations, when the graph is still rough and many candidate pairs could yield improvements, the filter is lenient.
In later iterations, when the graph has largely converged and most candidate pairs are wasteful (as shown in Figure~\ref{fig:waste}), the filter becomes stricter and skips more pairs.
The filter thus naturally tracks the convergence of the algorithm, applying more aggressive filtering precisely when it is most needed.

\subsection{False Negative Probability}
\label{sec:false_negative}

The filter may incorrectly skip a pair $(u_1, u_2)$ whose true distance would improve the graph---a false negative.
By the concentration bound (Equation~\ref{eq:concentration}), the probability that the projected distance of a true neighbor exceeds the threshold from a \emph{single} endpoint's perspective is at most $1 - p_\tau$.
Since we require both endpoints to agree (AND condition), the probability of a false negative for a pair that is a true neighbor of \emph{both} endpoints is at most:
\begin{equation}
    P_{\text{false negative}} \leq (1 - p_\tau)^2
    \label{eq:fn_prob}
\end{equation}

\begin{table}[t]
\centering
\caption{False negative probability bounds for different confidence levels.}
\label{tab:fn_prob}
\begin{tabular}{ccc}
\toprule
$p_\tau$ & $1 - p_\tau$ & $(1 - p_\tau)^2$ \\
\midrule
0.99 & 1.00\% & 0.01\% \\
0.95 & 5.00\% & 0.25\% \\
0.90 & 10.0\% & 1.00\% \\
0.85 & 15.0\% & 2.25\% \\
0.80 & 20.0\% & 4.00\% \\
\bottomrule
\end{tabular}
\end{table}

Table~\ref{tab:fn_prob} summarizes these bounds.
Note that these are \emph{worst-case} bounds---the actual false negative rate in practice depends on the distribution of projected distances and is often lower than the theoretical bound.
Furthermore, a false negative in one iteration does not necessarily result in a permanent recall loss: the same pair may be regenerated as a candidate in a subsequent iteration, at which point the tighter $d_k$ threshold may allow it to pass the filter.

\subsection{Comparison of Filtering Approaches}
\label{sec:filter_comparison}

Table~\ref{tab:filter_comparison} summarizes the key differences between the two filtering strategies.

\begin{table}[t]
\centering
\caption{Comparison of collision-count and projection-based filtering.}
\label{tab:filter_comparison}
\begin{tabular}{lcc}
\toprule
\textbf{Property} & \textbf{Collision} & \textbf{Projection} \\
\midrule
Signal type & Discrete (bits) & Continuous ($\mathbb{R}^m$) \\
Per-pair cost & $O(1)$ & $O(m)$ \\
Theoretical guarantee & None & Chi-squared bound \\
Threshold & Heuristic ($\mu$) & Derived ($\chi^2_{p_\tau}$) \\
Adaptive to convergence & Indirect & Direct via $d_k(u)$ \\
Preprocessing & $O(n \cdot H)$ & $O(n \cdot m \cdot d)$ \\
Memory & $O(n)$ (128 bits each) & $O(n \cdot m)$ \\
Sensitivity to $d$ & Low & Favorable for large $d$ \\
\bottomrule
\end{tabular}
\end{table}

The collision-count filter has the advantage of a lower per-pair cost ($O(1)$ vs.\ $O(m)$), but this advantage is outweighed by the projection filter's theoretically grounded threshold, continuous distance signal, and natural adaptivity to graph convergence.
The experimental comparison in Section~\ref{sec:experiments} confirms that the projection-based filter delivers substantially better results in practice.
